% Lösungen Einheit 2 von 2
\einheitenkopf[][2 Diagramme beurteilen und Boxplot]{Einheit 2 von 2 · Diagramme beurteilen und Boxplot}

\erg{25}{a) beginnt bei 0 \quad b) beginnt nicht bei 0, Startwert 60 \quad c) beginnt nicht bei 0, Startwert 1000 \quad d) beginnt bei 0}
\erg{26}{a) ja \quad b) nein (40 Tausend) \quad c) ja (44 Tausend < 46 Tausend) \quad d) nein (50 Tausend = 50\,000 Besucher) \quad e) ja (40 → 36 Tausend) \quad f) nein: von 2021 bis 2022 stieg die Zahl (44 → 46 Tausend)}
\erg{27}{a) falsch: die Hälfte von 180 ist 90, Sport hat 76 \quad b) wahr: 30 + 30 : 3 = 40 \quad c) wahr: 180 : 15 = 12 \quad d) falsch: 12 : 180 $\approx$ 6,7\,\% \quad e) nicht entscheidbar: das Diagramm unterscheidet nicht nach Mädchen und Jungen \quad f) falsch: Sport ist am häufigsten, aber Deutsch hat 22, die Hälfte von 40 wäre 20}
\erg{28}{a) Nein: 2023 waren es 420, 2024 440 Mitglieder. Die Achse beginnt bei 400, die Säulen zeigen nur 20 und 40. \quad b) 20 : 420 $\approx$ 4,8\,\% \quad c) Säulen bis 420 und 440, fast gleich hoch: der Zuwachs ist klein.}
\erg{29}{a) 1,20 € und 1,60 € \quad b) In Diagramm B: seine Achse ist gedehnt (dieselbe Höhe reicht nur bis 1,60 € statt bis 6 €), dieselben Werte wirken steiler. \quad c) 0,40 : 1,20 $\approx$ 33,3\,\%}
\erg{30}{a) Der Verkauf sinkt jedes Jahr, bis 2023 um 80 DVDs im Jahr, 2024 nur noch um 40. \quad b) Nicht haltbar: mit −80 im Jahr wären es 2026 genau 0, zuletzt sank der Verkauf aber nur um 40 (dann 2026 noch 80 DVDs); über die Zeit nach 2024 sagt das Diagramm nichts.}
\erg{31}{a) Die Aussage ist falsch: von 2020 bis 2021 stieg die Zahl (112 → 116 Mio.). Die Achse beginnt bei 90 Mio., die Säulen zeigen nur den Teil darüber: 2019 wirkt mehr als fünfmal so hoch wie 2024 (32 gegen 6 Mio. über 90), tatsächlich sank die Zahl von 122 auf 96 Mio., um etwa 21\,\%.}
\erg{32}{a) 260 cm \quad b) 450 cm \quad c) 340 cm \quad d) 310 cm und 380 cm \quad e) 450 cm − 260 cm = 190 cm \quad f) 75\,\%}
\erg{33}{a) Box von 6 bis 12, Strich bei 9, Antennen bis 3 und 18 \quad b) Minimum 4, unteres Quartil 6,5, Median 10, oberes Quartil 13, Maximum 17 \quad c) geordnet 5, 7, 8, 9, 10, 11, 12, 13, 14, 15, 16, 19: Minimum 5, unteres Quartil 8,5, Median 11,5, oberes Quartil 14,5, Maximum 19}
\erg{34}{a) 7a (13 gegen 12) \quad b) 7a (13 gegen 7) \quad c) 7b (Box 3 Punkte breit gegen 5) \quad d) Nur zum Teil: Das beste Ergebnis allein sagt wenig. Der Median der 7a ist höher (13 gegen 12), aber die 7a streut stärker: ihr unteres Quartil (10) und ihr Minimum (6) liegen unter denen der 7b (11 und 9).}
\erg{35}{a) Fehler: die Säulenhöhe über dem Achsenanfang (5 und 15) für den Wert genommen. Richtig: B hat 95, A 85 verkaufte Eis, also 10 mehr, etwa 11,8\,\% mehr. \quad b) Fehler: die Box mit der Spannweite verwechselt; die Box zeigt die mittlere Hälfte der Werte. Richtig: Spannweite 41 − 12 = 29.}
\erg{36}{a) Die Höhe einer Säule soll den Wert zeigen; beginnt die Achse nicht bei 0, stimmen die Höhenverhältnisse nicht mehr mit den Werten überein. \quad b) Nein: Temperaturen unter 36\,°C kommen hier nicht vor, und man liest die Werte an der Skala ab, nicht an Säulenhöhen; wichtig ist, dass die Achse gut beschriftet ist.}
\erg{37}{a) A: 960 → 990 Tausend, Zuwachs 30 Tausend, 30 : 960 $\approx$ 3,1\,\% \quad B: 400 → 440 Tausend, Zuwachs 40 Tausend, 40 : 400 = 10\,\% \quad b) Übertrieben: die Achse von A beginnt bei 950, die Säule 2024 wirkt viermal so hoch (40 gegen 10 über 950), der Zuwachs beträgt nur etwa 3,1\,\%. B ist in Tausend und in Prozent stärker gewachsen.}
